\section{Limitations}

Several limitations constrain the interpretation of HSAP results:

\textbf{Parameter values are illustrative.} The endocrine update equations, behavioral output functions, and sink mechanism are hypothesized, not empirically derived. Species-specific calibration is future work.

\textbf{Endocrine variables are proxies.} Testosterone, cortisol, and estrogen in the model are continuous values on $[0, 1]$ that represent qualitative endocrine states, not calibrated physiological measurements.

\textbf{Population trajectories are partly non-identifiable.} This is both a result (Section~\ref{sec:null}) and a limitation: HSAP cannot be tested using population data alone.

\textbf{Factorial scenarios use 30 seeds, not 50.} The mixed-depth design reflects a deliberate trade-off between statistical power for core scenarios and parameter-space exploration for factorial scenarios. If reviewers require equal depth, Set3 can be rerun at 50 seeds.

\textbf{HSAP does not model human sexual orientation.} Non-reproductive sexual or social behavior is treated only as species-specific ethological behavior and requires explicit supporting data. Species-specific calibration is future work.

\textbf{Well-mixed population.} The model assumes a well-mixed population without spatial structure or dispersal. This is a deliberate simplification to isolate endocrine-behavior mechanisms. Spatial extension is a natural next step.

\textbf{Linear endocrine model.} The endocrine update equations are linear combinations of environmental and social factors. Real endocrine systems involve nonlinear feedback, receptor dynamics, and individual genetic variation. The linear model captures first-order effects only.

%% TODO: Consider adding parameter identifiability analysis
%% TODO: Consider adding discussion of model selection criteria
